\subsection{The actual period-product germ}

\begin{result}{A.4.39}{Two normalized jets of the period product}{Proved}
Let an algebraic real short model have three distinct real roots
$e_1>e_2>e_3$, and put
\[
 \delta=e_1-e_3,\qquad m_0=\frac{e_2-e_3}{\delta}\in(0,1),
 \qquad F(m)=K(m)K(1-m),
\]
where $K$ is the complete elliptic integral in the parameter
convention. At $m=m_0$ the actual coefficients satisfy
\[
 h=\delta\left(m(1-m)\frac{F'}F-\frac{2m-1}{3}\right),
 \qquad c=\frac{\delta\pi}{4F}.
\]
The two numbers $F'(m_0)/F(m_0)$ and $F''(m_0)/F(m_0)$ are never
both algebraic. Consequently no nonzero constant normalization of
the actual germ $F(m_0+z)$ has all Taylor coefficients algebraic.
In particular it is not an ordinary $G$-function germ over $\Qb$.
This applies to the fixed non-CM $E_*$ and to the CM case $m_0=1/2$.

For comparison, there are explicit transcendental $G$-functions in
rational Fuchsian systems whose logarithmic derivative is algebraic
at an ordinary nonzero algebraic point.
\end{result}

\begin{proof}
Write $K=K(m)$, $L=K(1-m)$, $E=E(m)$ and $J=E(1-m)$.
The full-period formulas with the sign convention of A.4.1 are
\[
 \omega_1=\frac{2K}{\sqrt\delta},\quad
 \omega_2=\frac{2iL}{\sqrt\delta},\quad
 \frac{\eta_1}{\omega_1}
       =\delta\left(\frac EK-\frac{2-m}{3}\right),\quad
 \frac{\eta_2}{\omega_2}
       =\delta\left(\frac{1+m}{3}-\frac JL\right).
\]
The last two ratios are $h+c$ and $h-c$.
Differentiation of the defining integrals and integration of the
zero boundary term
$d(\sin\theta\cos\theta/\sqrt{1-m\sin^2\theta})$ give
\[
 K'=\frac{E-(1-m)K}{2m(1-m)},\qquad
 L'=-\frac{J-mL}{2m(1-m)}.
\]
Thus
\[
 \frac{F'}F=\frac{E/K-J/L+2m-1}{2m(1-m)},
\]
which proves the asserted formula for $h$.
Legendre's relation is $EL+JK-KL=\pi/2$, and gives both
$c=\delta\pi/(4F)$ and
\[
 W:=KL'-K'L=-\frac{\pi}{4m(1-m)}.
\]
These derivatives concern the family with fixed $\delta$ and roots
$(2-m)\delta/3$, $(2m-1)\delta/3$, $-(1+m)\delta/3$;
at $m_0$ it is exactly the given curve.

Both $K$ and $L$ satisfy $Y''+pY'+qY=0$, with
\[
 p=\frac{1-2m}{m(1-m)},\qquad q=-\frac1{4m(1-m)}.
\]
Expanding $F'=K'L+KL'$ and
$F''=-pF'-2qF+2K'L'$ proves
\[
 (F')^2-2FF''-2pFF'-4qF^2=W^2.
\]
With $r=F'/F$ and $s=F''/F$, this is the exact jet identity
\begin{equation}\label{ell:period-product-jets}
 s=\frac{r^2-2pr-4q}{2}
       -\frac{c^2}{2\delta^2m^2(1-m)^2}.
\end{equation}
If $r,s$ were both algebraic, the first formula of the result would
make $h$ algebraic, and \eqref{ell:period-product-jets} would make
$c^2$, hence $c$, algebraic. This contradicts the joint exclusion
proved in A.4.33. Therefore $r,s$ are never both algebraic.
If a nonzero scalar multiple of $F(m_0+z)$ had all coefficients
algebraic, its nonzero constant coefficient and its first two
coefficients would make $r,s$ algebraic, giving the same contradiction.
At $m_0=1/2$ one has $h=r=0$ and
$s=2-8c^2/\delta^2$; hence the obstruction persists in this proved
CM example. In the attempted algebraic-$h$ case, $r$ is algebraic
and \eqref{ell:period-product-jets} identifies the transcendental
second jet explicitly. No statement here decides whether $r$ is
algebraic at $E_*$.

The product also has a logarithmic germ at the original singular
center: $K(m)\to\pi/2$ and
$K(1-m)\sim-\tfrac12\log m$ as $m\to0^+$.
The latter follows by comparing the endpoint part of the defining
integral with $\int dt/\sqrt{t^2+m}$.
Thus this singular-center product is not an ordinary power series
either. Its actual scalar differential equation, obtained by
multiplication of the two second-order equations, is
\[
 F'''+3pF''+(p'+4q+2p^2)F'+(2q'+4pq)F=0.
\]
Having this rational differential equation does not change the
Taylor-coefficient obstruction just proved.

For the last assertion, for every positive integer $M$ take
\[
 G_M(z)=1+2\log(1-2Mz)-\log(1-3Mz),\qquad
 R_M(z)=-\frac{4M}{1-2Mz}+\frac{3M}{1-3Mz}.
\]
Its degree-$n$ Taylor coefficient, for $n\ge1$, is
$(3^n-2^{n+1})M^n/n$. The conjugate-size and common-denominator
bounds are exponential, and $(1,G_M)^{\mathsf T}$ solves the
rational Fuchsian system with matrix
$\left(\begin{smallmatrix}0&0\\R_M&0\end{smallmatrix}\right)$.
The derivative $R_M$ has residues $2,-1$, giving nonzero additive
logarithmic monodromy. Finite branching of algebraic functions
therefore proves that $G_M$ is transcendental and $1,G_M$ are
linearly independent over $\mathbf C(z)$.
At $\alpha=1/(4M)$, both logarithm arguments are positive,
$(1-2M\alpha)^2=1-3M\alpha$, and
\[
 G_M(\alpha)=1,\qquad G_M'(\alpha)=4M.
\]
This point is ordinary and strictly inside the radius $1/(3M)$.
It is the required exact pointwise example.
\end{proof}

The certificate \texttt{check\_period\_product\_gfunction\_route.py}
checks the derivative, Wronskian, symmetric-square and counterexample
identities. The small-rational-point theorem in
\cite[Theorem I, original p.7263]{EllChudnovsky1984PNAS} retains its
rational-coefficient and quantitative smallness hypotheses;
A.4.39 supplies no eligible ordinary $G$-function germ by normalizing
the actual period product. The infinite-family dimension theorem
\cite[Theorem 3, p.4 of the opened manuscript]{EllFischlerRivoal2018GValues}
also does not identify its iterated-integral input with these jets.
The demonstrated failures concern these specified substitutions.
The main nonzero algebraic-$h$ question remains unchanged.
